\usepackage[T2A]{fontenc}
\usepackage[utf8]{inputenc}
\usepackage[russian]{babel}

% ============================================================================
% Макет по образцу серии «МЗ Пресс — МФТИ» (Натан, Горбачёв, Гуз, 2007 и др.):
% A5, Times, 11 pt, одинарный интервал, абзац 10 mm, поля ~1,6–1,8 cm.
% ============================================================================

\usepackage{tempora}

\usepackage{amsmath}
\usepackage{newtxmath}  % не подключать amssymb/amsfonts --- дублируют \Bbbk и др.

\usepackage[
  a5paper,
  twoside,
  inner=1.8cm,
  outer=1.8cm,
  top=1.7cm,
  bottom=2.3cm,
  bindingoffset=3mm,
  headheight=12pt,
  headsep=8pt,
  footskip=20pt
]{geometry}

\usepackage{setspace}
\singlespacing

\setlength{\parindent}{1.0cm}
\setlength{\parskip}{0pt plus 0pt minus 0pt}

\usepackage{fancyhdr}
\pagestyle{fancy}
\fancyhf{}
\fancyfoot[C]{\normalsize\thepage}
\renewcommand{\headrulewidth}{0pt}
\fancypagestyle{plain}{%
  \fancyhf{}%
  \fancyfoot[C]{\normalsize\thepage}%
  \renewcommand{\headrulewidth}{0pt}%
}

\usepackage{microtype}
\usepackage{array}
\usepackage{tabularx}
\usepackage{ragged2e}
\usepackage{indentfirst}
\usepackage{natbib}
\usepackage{graphicx}
\usepackage{grffile}
\usepackage{subcaption}
\usepackage{booktabs}
\usepackage{longtable}
\usepackage{hyperref}
\usepackage{xcolor}
\usepackage{tcolorbox}
\tcbuselibrary{skins,breakable,listings}

\hypersetup{colorlinks=true,linkcolor=blue!60!black,citecolor=blue!60!black,urlcolor=blue!60!black}

% Арабская нумерация с титула; без сброса счётчика при \mainmatter.
\makeatletter
\renewcommand\frontmatter{%
  \if@twoside
    \cleardoublepage
  \else
    \clearpage
  \fi
  \@mainmatterfalse
  \pagenumbering{arabic}%
}
\renewcommand\mainmatter{%
  \if@twoside
    \cleardoublepage
  \else
    \clearpage
  \fi
  \@mainmattertrue
}
\makeatother

% Оглавление: главы и \section, без \subsection и \subsubsection.
\setcounter{tocdepth}{1}

\DeclareMathOperator*{\plim}{\mathop{plim}}
\DeclareMathOperator{\prob}{\mathbf{P}\!}
\DeclareMathOperator{\cov}{cov}
\DeclareMathOperator{\corr}{corr}
\DeclareMathOperator{\med}{med}
\DeclareMathOperator{\rank}{rank}
\DeclareMathOperator{\arctanh}{arctanh}
\DeclareMathOperator*{\argmin}{arg\,min}
\DeclareMathOperator*{\argmax}{arg\,max}
\DeclareMathOperator{\diag}{diag}

\DeclareMathOperator{\logit}{logit}
\DeclareMathOperator{\FWER}{FWER}
\DeclareMathOperator{\FDR}{FDR}
\DeclareMathOperator{\FDP}{FDP}

\newtheorem{definition}{Определение}[chapter]
\newtheorem{remark}{Замечание}[chapter]

\newtcblisting{codefragment}{
  listing only,
  listing options={
    basicstyle=\small\ttfamily,
    breaklines=true,
    language=Python,
  },
  colback=gray!10,colframe=gray!50,breakable
}

\newcommand{\FigDir}{figures}
\newcommand{\FigWidth}{0.92\linewidth}  % полная ширина иллюстрации на A5

% Репозиторий курса ПСАД (данные лабораторных работ).
\newcommand{\PSADRepo}{\url{https://github.com/andriygav/PSAD}}
\newcommand{\PSADPath}[1]{\url{https://github.com/andriygav/PSAD/tree/master/#1}}

% Карточка критерия (формат слайдов лекции: выборка, H_0, T, null, p).
\newenvironment{critcard}[1]{%
  \par\medskip\noindent\textbf{#1}\par\smallskip
  \begingroup\footnotesize
  \setlength{\tabcolsep}{4pt}%
  \begin{tabular}{@{}>{\raggedright\arraybackslash}p{0.24\linewidth}>{\raggedright\arraybackslash}p{0.71\linewidth}@{}}%
}{%
  \end{tabular}\endgroup\par\medskip
}

% Таблица обозначений: обычный tabular (tabularx ломается при \begin/\end в макросах).
\newcommand{\NotationTableBegin}{%
  \par\addvspace{0.35\baselineskip}%
  \footnotesize
  \setlength{\tabcolsep}{3pt}%
  \setlength{\extrarowheight}{1.5pt}%
  \begin{tabular}{@{}>{\raggedright\arraybackslash}p{0.26\linewidth}>{\raggedright\arraybackslash}p{0.70\linewidth}@{}}%
}
\newcommand{\NotationTableEnd}{%
  \end{tabular}%
  \normalsize
  \par\addvspace{0.25\baselineskip}%
}
